% From mitthesis package
% Version: 1.01, 2023/06/19
% Documentation: https://ctan.org/pkg/mitthesis
%
% If you don't want an initial indentation, do \noindent at the start of the abstract

This research project is an exploration of the use of \ac{AR} for \ac{CL}. Current literature shows a growing interest in the use of interactive technologies for teaching and the overall
impact of \ac{TEL} projects experimenting with the implementation of \ac{AR}. This technology offers interesting visualisation tools to merge digital information with the physical context
of the user, creating unique interactive experiences and boosting engagement for the students, but it had been tested mostly on single-user settings.

The proposal of this thesis is that the unique visualisation capabilities of \ac{AR} and the connectivity offered by the mobile hardware often used for its deployment can provide an
advantage for effective \ac{CL}. For this purpose, \emph{WorkshopAR} was developed following a user-centred design and implemented within the context of a project-based course in
construction design. This application offers the students different tools to help them structure and manage their collaborative work.

The Industry Project course was used as a case study to observe the impact that the introduction of \emph{WorkshopAR} produced in the behaviour of the students during their classroom
activities. The main objective was to promote positive changes through the application, helping the students develop and apply important skills for effective collaboration and for a
successful development of their project. The risks and uncertainties of such a strong intervention with technology in social, collaborative activities were also taken into consideration,
and the observation process also identified negative or undesirable changes in behaviour that could be acknowledged through the iterative development process used.

The case study also offered an opportunity to document the design and implementation of the software. The development process integrated the functional requirements of the tool with the
learning goals of the course and tied the iterative development with what was observed in the classroom. This allowed a flexible and adaptable design that boosted the interesting
approaches that students and teaching staff were proposing while also helped in redesign those elements that were promoting problematic behaviours.

The results obtained showed that students benefited the most with a strong sense of structure when initially establishing relationships with the group, and appreciated more flexible
activities once the work dynamics were clearly defined and producing results. The visual and interactive elements proposed by \emph{WorkshopAR} helped the students understand and be aware
of the need of those structures, promoting discussions and guiding work using more impactful collaborative approaches, but failed to integrated in a more substantial way in the workflows
that students had already created using a variety of other tools that were better integrated into each other. \emph{WorkshopAR} was able to pivot when aspects of its design were promoting
behaviours that were not in line with the learning goals of the course or the research, and several opportunities in the creation of collaborative spaces and communication mediated by \ac{AR}
can be worked upon in future research.